% ============================================================
% 第7章 实验结果与关键发现
% 【负责人：建模线（7.1 基线结果已预填实测数字；7.2/7.3 待补充）】
% 赛题要求“关键发现”为说明文档必备内容。
% ============================================================
\section{实验结果与关键发现}

\subsection{基线结果（已完成）}

\begin{table}[H]
\centering
\caption{逻辑回归基线五折折外结果（2.0 版数据）}
\label{tab:baseline-res}
\small
\begin{tabular}{lccccc}
\toprule
折号 & 0 & 1 & 2 & 3 & 4 \\
\midrule
样本数 & 101 & 100 & 100 & 100 & 99 \\
正类数 & 26 & 25 & 25 & 25 & 25 \\
折内 AUC & 0.828 & 0.812 & 0.792 & 0.716 & 0.750 \\
\bottomrule
\end{tabular}

\medskip
\begin{tabularx}{\textwidth}{lL}
\toprule
整体指标 & 数值 \\
\midrule
折外 AUC（OOF） & \textbf{0.7766}（95\% CI：0.726--0.826，车辆级配对 bootstrap） \\
average\_precision & 0.581 \\
brier / log\_loss & 0.145 / 0.456（常数先验 brier = 0.189） \\
Recall@50 / Precision@50 & 0.286 / 0.720（lift = 2.86） \\
\bottomrule
\end{tabularx}
\end{table}

\begin{TODO}
请建模线补充基线结果的解读（各特征系数方向与符号、折 3 偏低的原因分析等）。
\end{TODO}

\subsection{最终模型结果}

\begin{TODO}
负责人：建模线。需包含：(1) 主表——各候选模型在统一折分下的折外 AUC 对比
（固定特征比较模型，固定模型比较特征）；(2) 方案 A/B（单一 vs 集成）对比结论；
(3) 消融实验（清洗版本 v2 vs v4、特征模块逐项加入的贡献）；(4) 最终提交模型的
选择依据与 500 台车全量预测的生成方式。
\end{TODO}

\subsection{关键发现}

\begin{TODO}
负责人：全员按线索提供、建模线汇总。已有可写素材：
(1) 清洗线——(0,0) 占位坐标事故若按坐标过滤将丢失真实正类（3.4 节）；
坐标无效不影响标签守恒；7/31 越界数据整日隔离的影响评估；
(2) 数据线——2.0 版换版导致的名单与正类变化（126 vs 149），任何 1.0 结论
不可沿用；
(3) 建模线——待补：哪些行为特征最有区分力、IMU 特征的增益、失败案例复盘。
\end{TODO}
